% =====================================================================
\chapter{Números complejos en forma binómica}
\label{tema:5}
% =====================================================================

Los números complejos amplían los reales con una raíz cuadrada de $-1$. Con ellos toda ecuación polinómica tiene solución, y las operaciones de giro y homotecia del plano se escriben como productos. Este tema construye $\C$, demuestra las propiedades del conjugado y del módulo y resuelve ecuaciones de segundo grado con coeficientes complejos. Se sigue la §2.4 de~\cite{osat}, la sección~5 de las notas de 18.03 del MIT~\cite{mit1803} y la §1.1 de las notas de «Vectors and Matrices» de Cambridge~\cite{cambvm}.

\section{El cuerpo de los números complejos}

\begin{definicion}[Números complejos]\label{def:complejos}
Un \emph{número complejo} es una expresión $z=a+bi$ con $a,b\in\R$; se llaman \emph{parte real} $\Real z=a$ y \emph{parte imaginaria} $\Imag z=b$. Dos complejos son iguales si tienen iguales las partes reales y las imaginarias. Con $z=a+bi$ y $w=c+di$ se definen
\[
z+w=(a+c)+(b+d)i,\qquad zw=(ac-bd)+(ad+bc)i .
\]
El conjunto de los números complejos se denota $\C$. Se identifica el real $a$ con $a+0i$, y se escribe $i=0+1i$.
\end{definicion}

La regla del producto es la que resulta de operar como con polinomios en $i$ y sustituir $i^{2}=-1$: en efecto, con la definición, $i\cdot i=(0\cdot0-1\cdot1)+(0\cdot1+1\cdot0)i=-1$. Geométricamente, $a+bi$ se representa como el punto $(a,b)$ del \emph{plano complejo}, con el eje real horizontal y el imaginario vertical (\cref{fig:argand}).

\begin{teorema}\label{teo:cuerpo}
Con estas operaciones, $\C$ es un cuerpo: la suma y el producto son asociativos y conmutativos, el producto es distributivo respecto de la suma, $0$ y $1$ son los elementos neutros, y todo $z=a+bi\ne0$ tiene inverso
\[
z^{-1}=\frac{a}{a^{2}+b^{2}}-\frac{b}{a^{2}+b^{2}}\,i .
\]
\end{teorema}

\begin{proof}
Las propiedades de la suma son las de $\R^{2}$ componente a componente. La conmutatividad del producto se ve en la fórmula, simétrica en $(a,b)$ y $(c,d)$; la asociatividad y la distributividad se comprueban desarrollando ambos miembros (el cálculo, rutinario, se propone en~\cite[§5.1]{mit1803}). Para el inverso, si $z\ne0$ entonces $a^{2}+b^{2}>0$ y
\[
(a+bi)\Bigl(\frac{a}{a^{2}+b^{2}}-\frac{b}{a^{2}+b^{2}}i\Bigr)=\frac{a^{2}+b^{2}}{a^{2}+b^{2}}+\frac{-ab+ab}{a^{2}+b^{2}}\,i=1.\qedhere
\]
\end{proof}

\begin{ejemplo}[Potencias de $i$]\label{ej:potencias-i}
$i^{1}=i$, $i^{2}=-1$, $i^{3}=-i$, $i^{4}=1$, y a partir de ahí se repiten: $i^{n}$ solo depende del resto de dividir $n$ entre~4. Por ejemplo, $i^{103}=i^{4\cdot25+3}=i^{3}=-i$.
\end{ejemplo}

\section{El conjugado}

\begin{definicion}[Conjugado]
El \emph{conjugado} de $z=a+bi$ es $\bar z=a-bi$. Geométricamente, es la reflexión de $z$ en el eje real.
\end{definicion}

\begin{proposicion}[Propiedades del conjugado]\label{prop:conjugado}
Para todos $z,w\in\C$:
\begin{enumerate}
  \item $\overline{z+w}=\bar z+\bar w$, \quad $\overline{zw}=\bar z\,\bar w$, \quad $\overline{z/w}=\bar z/\bar w$ si $w\ne0$, \quad $\bar{\bar z}=z$;
  \item $z+\bar z=2\Real z$ y $z-\bar z=2i\Imag z$;
  \item $z$ es real si y solo si $z=\bar z$;
  \item $z\bar z=(\Real z)^{2}+(\Imag z)^{2}$, que es un número real mayor o igual que~0.
\end{enumerate}
\end{proposicion}

\begin{proof}
Sean $z=a+bi$ y $w=c+di$. (i) $\overline{z+w}=(a+c)-(b+d)i=\bar z+\bar w$. Para el producto, $\overline{zw}=(ac-bd)-(ad+bc)i$, y $\bar z\bar w=(a-bi)(c-di)=(ac-bd)+(-ad-bc)i$, que es lo mismo. Para el cociente, $\bar w\,\overline{(z/w)}=\overline{w\cdot(z/w)}=\bar z$. (ii) y (iii) son inmediatas. (iv) $(a+bi)(a-bi)=a^{2}-(bi)^{2}=a^{2}+b^{2}$.
\end{proof}

La propiedad~(iv) da el método práctico para dividir: se multiplican numerador y denominador por el conjugado del denominador,
\[
\frac{z}{w}=\frac{z\,\bar w}{w\,\bar w}=\frac{z\,\bar w}{c^{2}+d^{2}} .
\]

\begin{ejemplo}\label{ej:division}
$\dfrac{3+i}{1-2i}=\dfrac{(3+i)(1+2i)}{1+4}=\dfrac{3+6i+i+2i^{2}}{5}=\dfrac{1+7i}{5}=\dfrac15+\dfrac75i$.
\end{ejemplo}

\section{El módulo}

\begin{definicion}[Módulo]
El \emph{módulo} de $z=a+bi$ es $\abs{z}=\sqrt{a^{2}+b^{2}}=\sqrt{z\bar z}$, la distancia de $z$ al origen. La distancia entre $z$ y $w$ es $\abs{z-w}$.
\end{definicion}

\begin{proposicion}[Propiedades del módulo]\label{prop:modulo}
Para todos $z,w\in\C$:
\begin{enumerate}
  \item $\abs{z}\ge0$, con igualdad si y solo si $z=0$; \quad $\abs{\bar z}=\abs{z}$;
  \item $\abs{\Real z}\le\abs{z}$ y $\abs{\Imag z}\le\abs{z}$;
  \item $\abs{zw}=\abs{z}\abs{w}$ y, si $w\ne0$, $\abs{z/w}=\abs{z}/\abs{w}$;
  \item (desigualdad triangular) $\abs{z+w}\le\abs{z}+\abs{w}$;
  \item $\bigl|\abs{z}-\abs{w}\bigr|\le\abs{z-w}$.
\end{enumerate}
\end{proposicion}

\begin{proof}
(i) y (ii) se leen en la fórmula: $a^{2}\le a^{2}+b^{2}$. (iii) Por la \cref{prop:conjugado}, $\abs{zw}^{2}=zw\,\overline{zw}=z\bar z\,w\bar w=\abs{z}^{2}\abs{w}^{2}$, y se toman raíces no negativas; el cociente se sigue de aplicar lo anterior a $(z/w)\cdot w=z$. (iv) Por la \cref{prop:conjugado}(ii), $z\bar w+\bar zw=z\bar w+\overline{z\bar w}=2\Real(z\bar w)$, luego
\[
\abs{z+w}^{2}=(z+w)(\bar z+\bar w)=\abs{z}^{2}+\abs{w}^{2}+2\Real(z\bar w)\le\abs{z}^{2}+\abs{w}^{2}+2\abs{z\bar w}=(\abs{z}+\abs{w})^{2},
\]
usando (ii) y (iii). (v) Por (iv), $\abs{z}=\abs{(z-w)+w}\le\abs{z-w}+\abs{w}$, luego $\abs{z}-\abs{w}\le\abs{z-w}$; intercambiando los papeles, $\abs{w}-\abs{z}\le\abs{z-w}$.
\end{proof}

\section{Ecuaciones polinómicas}

\begin{proposicion}[Raíces cuadradas]\label{prop:raiz-cuadrada}
Sea $w=a+bi$ con $b\ne0$. Las soluciones de $z^{2}=w$ son $z=\pm(x+yi)$ con
\[
x=\sqrt{\frac{\abs{w}+a}{2}},\qquad y=\operatorname{signo}(b)\sqrt{\frac{\abs{w}-a}{2}} .
\]
\end{proposicion}

\begin{proof}
Con $z=x+yi$, la ecuación $z^{2}=w$ equivale a $x^{2}-y^{2}=a$ y $2xy=b$. Tomando módulos, $\abs{z}^{2}=\abs{w}$, es decir, $x^{2}+y^{2}=\abs{w}$. Sumando y restando con la primera ecuación, $x^{2}=\frac{\abs{w}+a}{2}$ e $y^{2}=\frac{\abs{w}-a}{2}$, ambos positivos porque $\abs{w}>\abs{a}$ cuando $b\ne0$. La condición $2xy=b$ fija que $x$ e $y$ tengan el producto del signo de $b$, lo que deja exactamente los dos pares del enunciado; y estos pares la cumplen, porque $(2xy)^{2}=4\cdot\frac{\abs{w}^{2}-a^{2}}{4}=b^{2}$.
\end{proof}

\begin{ejemplo}\label{ej:raiz}
Para $w=3+4i$, $\abs{w}=5$, $x=\sqrt{4}=2$, $y=\sqrt1=1$: las raíces son $\pm(2+i)$. Comprobación: $(2+i)^{2}=4+4i+i^{2}=3+4i$.
\end{ejemplo}

\begin{proposicion}[Fórmula cuadrática en $\C$]\label{prop:cuadratica}
Sean $a,b,c\in\C$ con $a\ne0$, y sea $\delta\in\C$ con $\delta^{2}=b^{2}-4ac$. Las soluciones de $az^{2}+bz+c=0$ son $z=\dfrac{-b\pm\delta}{2a}$.
\end{proposicion}

\begin{proof}
Multiplicando por $4a\ne0$ y completando el cuadrado, la ecuación equivale a $(2az+b)^{2}=b^{2}-4ac=\delta^{2}$, es decir, $(2az+b-\delta)(2az+b+\delta)=0$. Como $\C$ es un cuerpo, un producto es nulo si y solo si lo es algún factor, de donde $2az+b=\pm\delta$.
\end{proof}

\begin{proposicion}[Raíces conjugadas]\label{prop:conjugadas}
Si $p(z)=\sum_{k=0}^{n}c_kz^{k}$ tiene coeficientes reales, entonces $p(\bar z)=\overline{p(z)}$ para todo $z$. En particular, si $z$ es raíz de $p$, también lo es $\bar z$.
\end{proposicion}

\begin{proof}
Por la \cref{prop:conjugado}(i), aplicada repetidamente, $\overline{z^{k}}=\bar z^{k}$ y $\overline{\sum c_kz^{k}}=\sum\bar c_k\bar z^{k}=\sum c_k\bar z^{k}$, porque $\bar c_k=c_k$. Si $p(z)=0$, entonces $p(\bar z)=\bar0=0$.
\end{proof}

\begin{teorema}[Teorema fundamental del álgebra]\label{teo:fundamental}
Todo polinomio de grado $n\ge1$ con coeficientes complejos tiene al menos una raíz en $\C$. En consecuencia, se factoriza como $c\prod_{k=1}^{n}(z-z_k)$ y tiene exactamente $n$ raíces contadas con multiplicidad.
\end{teorema}

La demostración requiere herramientas de análisis y se omite; véanse~\cite[§5.3]{mit1803} y~\cite{needham}. La consecuencia se obtiene por inducción sobre $n$, dividiendo por $z-z_1$ cada vez que se encuentra una raíz.

\section{El plano complejo y el producto escalar}

\begin{figure}[ht]
\centering
\begin{tikzpicture}[scale=1.1,font=\small,>={Stealth[length=2mm]}]
  \draw[->] (-1.2,0) -- (4.2,0) node[right] {$\Real$};
  \draw[->] (0,-2.4) -- (0,2.8) node[above] {$\Imag$};
  \coordinate (z) at (3,2);
  \coordinate (zb) at (3,-2);
  \draw[thick,->] (0,0) -- (z);
  \draw[thick,->,dashed] (0,0) -- (zb);
  \draw[dotted] (z) -- (3,0) node[below right] {$a$};
  \draw[dotted] (z) -- (0,2) node[left] {$b$};
  \fill (z) circle (0.05) node[above right] {$z=a+bi$};
  \fill (zb) circle (0.05) node[below right] {$\bar z=a-bi$};
  \node at (1.25,1.25) [rotate=33.7] {$\abs{z}$};
\end{tikzpicture}
\caption{El plano complejo. El número $z=a+bi$ es el punto $(a,b)$; su módulo es la distancia al origen y su conjugado, la reflexión en el eje real.}
\label{fig:argand}
\end{figure}

Al identificar $\C$ con $\R^{2}$, la suma de complejos es la suma de vectores, y el producto por un real $\lambda$ es el producto de un vector por un escalar. El producto de dos complejos no tiene análogo en $\R^{2}$ y es el objeto del \cref{tema:6}. El conjugado permite recuperar el producto escalar y el área orientada:

\begin{proposicion}\label{prop:escalar-complejo}
Sean $z=a+bi$ y $w=c+di$, identificados con $u=(a,b)$ y $v=(c,d)$. Entonces
\[
\Real(z\bar w)=ac+bd=u\cdot v,\qquad \Imag(\bar z w)=ad-bc=\det\begin{pmatrix}a&c\\ b&d\end{pmatrix}.
\]
\end{proposicion}

\begin{proof}
$z\bar w=(a+bi)(c-di)=(ac+bd)+(bc-ad)i$ y $\bar zw=(a-bi)(c+di)=(ac+bd)+(ad-bc)i$.
\end{proof}

\begin{observacion}[Errores frecuentes]
(a) La parte imaginaria de $a+bi$ es el número real $b$, no $bi$. (b) $\abs{z}^{2}=z\bar z$, pero en general $\abs{z}^{2}\ne z^{2}$: para $z=i$, $\abs{z}^{2}=1$ y $z^{2}=-1$. (c) La regla $\sqrt{zw}=\sqrt z\sqrt w$ falla en $\C$: con $z=w=-1$ daría $1=\sqrt1=\sqrt{-1}\sqrt{-1}=i^{2}=-1$. Todo número complejo no nulo tiene dos raíces cuadradas y no hay una elección de signo compatible con el producto. (d) En $\C$ no hay un orden compatible con las operaciones (\cref{eje:5-orden}): no tiene sentido escribir $z<w$ entre complejos no reales.
\end{observacion}

\begin{observacion}[Conexión con el curso]
Una matriz real puede tener autovalores complejos, y por la \cref{prop:conjugadas} aparecen en pares conjugados, porque el polinomio característico tiene coeficientes reales (módulo~3, tema~3). La \cref{prop:escalar-complejo} es la identidad que usa la forma compleja de RoPE (módulo~14, tema~4): cada par de coordenadas $(q_{2j},q_{2j+1})$ de una consulta se trata como un número complejo, y el producto escalar de dos pares es $\Real(z\bar w)$.
\end{observacion}

\begin{lecturas}
\item OpenStax, \emph{Algebra and Trigonometry}~\cite{osat}, §2.4 «Complex Numbers».
\item MIT 18.03, \emph{Supplementary Notes}~\cite{mit1803}, §5 «The algebra of complex numbers» (5.1 a 5.3).
\item Cambridge, \emph{Vectors and Matrices}~\cite{cambvm}, §1.1 «Basic properties».
\end{lecturas}

\section*{Ejercicios del tema 5}
\addcontentsline{toc}{section}{Ejercicios del tema 5}

\begin{ejercicio}{1}{eje:5-partes}
Para $z_1=3-4i$, $z_2=-2i$, $z_3=5$ y $z_4=-1+i$, indique la parte real, la parte imaginaria, el conjugado y el módulo, y sitúe cada número en el plano complejo.
\end{ejercicio}

\begin{ejercicio}{1}{eje:5-vf}
Justifique si son verdaderas o falsas para todos $z,w\in\C$: (a) $\abs{z+w}=\abs{z}+\abs{w}$; (b) $z\bar z$ es real y no negativo; (c) $\overline{z^{2}}=\bar z^{\,2}$; (d) $\Imag(3+2i)=2i$.
\end{ejercicio}

\begin{ejercicio}{1}{eje:5-potencias-i}
Calcule $i^{7}$, $i^{100}$, $i^{-1}$ e $i^{2026}$, y la suma $1+i+i^{2}+\dots+i^{10}$.
\end{ejercicio}

\begin{ejercicio}{2}{eje:5-operaciones}
Escriba en forma binómica $(3+2i)(1-4i)$, $(2-i)^{2}$, $\dfrac{2-i}{1+3i}$ y $(1+i)^{4}$.
\fuente{adaptado de OpenStax, \emph{Algebra and Trigonometry}~\cite{osat}, §2.4, ejercicios de la sección sobre operaciones con complejos.}
\end{ejercicio}

\begin{ejercicio}{2}{eje:5-inversos}
Calcule $i^{-1}$, $(1+i)^{-1}$ y $\dfrac{1+i}{2-i}$, y exprese $\abs{z^{-1}}$ en función de $\abs{z}$.
\fuente{adaptado de MIT 18.03, \emph{Supplementary Notes}~\cite{mit1803}, ejercicio~5.2.3.}
\end{ejercicio}

\begin{ejercicio}{2}{eje:5-cuadraticas}
Encuentre las raíces de $x^{2}+x+1$ y de $x^{2}+x+2$, y resuelva $x^{2}+ix+1=0$.
\fuente{adaptado de MIT 18.03, \emph{Supplementary Notes}~\cite{mit1803}, ejercicio~5.2.4.}
\end{ejercicio}

\begin{ejercicio}{2}{eje:5-suma-tripos}
Sean $x,y\in\R$ tales que $x+iy=\sum_{a=0}^{200}i^{a}+\sum_{b=1}^{50}i^{b}$. Calcule $xy$.
\fuente{adaptado de Cambridge, \emph{Mathematical Tripos} Part IA, Paper~1, 2019~\cite{tripos2019}, pregunta~1C(a).}
\end{ejercicio}

\begin{ejercicio}{2}{eje:5-raices-cuadradas}
Calcule las raíces cuadradas de $-5+12i$ y de $-2i$ con la \cref{prop:raiz-cuadrada}, y compruébelas elevando al cuadrado.
\end{ejercicio}

\begin{ejercicio}{2}{eje:5-cuadratica-compleja}
Resuelva $z^{2}-(3+i)z+(2+2i)=0$ y compruebe que la suma y el producto de las raíces son los esperados.
\end{ejercicio}

\begin{ejercicio}{2}{eje:5-lugares}
Describa geométricamente y dibuje los conjuntos (a) $\{z:\abs{z-1}=\abs{z+i}\}$ y (b) $\{z:\abs{z+3i}=3\abs{z}\}$.
\fuente{el apartado (b) está adaptado de Cambridge, \emph{Mathematical Tripos} Part IA, Paper~1, 2013~\cite{tripos2013}, pregunta~1C(b).}
\end{ejercicio}

\begin{ejercicio}{2}{eje:5-escalar}
Sean $z=1+2i$ y $w=3+i$. Calcule $z\bar w$ y compruebe la \cref{prop:escalar-complejo}: $\Real(z\bar w)$ es el producto escalar de $(1,2)$ y $(3,1)$. Calcule también el ángulo entre ambos vectores.
\end{ejercicio}

\begin{ejercicio}{3}{eje:5-modulo-producto}
Demuestre directamente, sin usar el conjugado, que si $z=a+bi$ entonces $z\bar z=a^{2}+b^{2}$ y que $\abs{wz}=\abs{w}\abs{z}$ desarrollando ambos miembros en coordenadas. ¿Qué identidad entre números reales se obtiene?
\fuente{adaptado de MIT 18.03, \emph{Supplementary Notes}~\cite{mit1803}, ejercicios~5.2.1 y~5.2.2.}
\end{ejercicio}

\begin{ejercicio}{3}{eje:5-igualdad-triangular}
Demuestre que, para $z,w\ne0$, se cumple $\abs{z+w}=\abs{z}+\abs{w}$ si y solo si $w=\lambda z$ para algún real $\lambda>0$.
\end{ejercicio}

\begin{ejercicio}{3}{eje:5-orden}
Demuestre que no existe en $\C$ un orden total $<$ compatible con las operaciones, es decir, que cumpla: si $z<w$ entonces $z+u<w+u$ para todo $u$, y si $0<z$ y $0<w$ entonces $0<zw$. Indicación: en un orden así, todo cuadrado no nulo es positivo.
\end{ejercicio}

\begin{ejercicio}{3}{eje:5-blaschke}
Sean $z,w\in\C$ con $\abs{z}=1$ y $\bar wz\ne1$. Demuestre que $\left|\dfrac{z-w}{1-\bar wz}\right|=1$.
\end{ejercicio}
